\section{Discusión comparativa}

\subsection{Qué variantes ganaron y por qué}
La Tabla~\ref{tab:variantes} resume la configuración elegida en cada caso. Hay dos
patrones que se repiten en ambos datasets. Para el árbol de decisión, exigir un mínimo
de 10 observaciones por hoja fue siempre la mejor variante: un árbol sin límites
memoriza el entrenamiento, y limitar el tamaño de las hojas reduce esa varianza sin
eliminar las reglas útiles. Para KNN, aumentar $k$ a 15 no fue la mejor variante en
ningún dataset: en Diabetes bajó el F1 macro de 0.719 a 0.681, porque con 45 diabéticos
en train un vecindario de 15 casos queda casi siempre dominado por la clase mayoritaria;
en Potabilidad subió apenas 0.011, menos que el cambio a distancia Manhattan.

En cambio, otras decisiones dependen del dataset. Ponderar las clases fue la mejor
variante de XGBoost y AdaBoost en Potabilidad, donde los modelos tienden a predecir
``no potable''. En Diabetes no fue la mejor variante en ninguna familia y bajó el F1 de
DT (0.783 a 0.751) y RF (0.836 a 0.816): dar más peso a 45 diabéticos lleva a los
árboles a crear reglas para casos aislados. La selección de atributos mejoró cuatro
de cinco familias en Diabetes (en RF sólo 0.003, una diferencia despreciable) y empeoró las cinco en Potabilidad. La explicación es la
forma de la relación entre variables y clase: en Diabetes la señal está casi toda en la
glucosa y es monótona (más glucosa, más riesgo), así que un filtro univariado la
encuentra; en Potabilidad la señal está en los extremos y en la combinación de sulfato
y \texttt{ph}, y el filtro descarta el \texttt{ph} aunque sea la segunda variable más
importante para el mejor modelo.

\begin{table}[H]\centering\small
\caption{Configuración elegida por validación cruzada en cada familia y dataset.}
\label{tab:variantes}
\begin{tabular}{lll}\toprule
Familia & Potabilidad & Diabetes \\ \midrule
KNN & Manhattan ($p=1$) & estándar + selección de 11 atributos \\
DT & mínimo 10 por hoja & mínimo 10 por hoja + selección \\
RF & 500 árboles & estándar + selección \\
XGBoost & clases balanceadas & tasa de aprendizaje 0.05 \\
AdaBoost & clases balanceadas & tasa 0.5 + selección \\ \bottomrule
\end{tabular}
\end{table}

\subsection{Sobreajuste y estabilidad}
Las Figuras~\ref{fig:resumen-potabilidad} y~\ref{fig:resumen-diabetes} comparan el F1 de
entrenamiento, validación cruzada y prueba. Random Forest y XGBoost llegan a F1 de
entrenamiento de 1.0 o casi 1.0 en ambos datasets: memorizan el entrenamiento. Aun así,
RF es el que mejor generaliza en Potabilidad, porque el promedio de cientos de árboles
cancela buena parte de sus errores individuales. AdaBoost con stumps es el extremo
opuesto en Potabilidad: casi no sobreajusta (0.575 en entrenamiento y 0.582 en prueba),
pero subajusta. KNN y DT quedan en medio.

La validación cruzada también muestra cuánto se puede confiar en cada diferencia. En
Potabilidad (unos 300 casos por pliegue) las desviaciones de los ganadores están entre
0.02 y 0.05, y la ventaja de RF sobre KNN, DT y AdaBoost es clara; frente a XGBoost
(0.639) la diferencia es menor que la desviación. En Diabetes (unos 58 casos y 9
diabéticos por pliegue) las desviaciones llegan a 0.12 y las cuatro familias de árboles
son estadísticamente indistinguibles. Una sola partición de validación, como la que se
usaba antes en el framework, habría dado un ganador sin advertir esa incertidumbre.

\subsection{Exactitud frente a F1 macro}
En Diabetes la exactitud de prueba (88.8--89.8\,\%) parece alta, pero la referencia que
siempre predice ``no diabetes'' ya obtiene 84.7\,\%. El F1 macro (0.75--0.78 frente a
0.459) y el recall de diabetes (0.47--0.53) muestran mejor lo que hace el modelo: detecta
la mitad de los diabéticos. En Potabilidad, AdaBoost estándar tiene una exactitud de
validación de 59.7\,\%, igual a la proporción de la clase mayoritaria: sin el F1 macro
se habría pasado por alto que el modelo casi no aprendió nada.

\subsection{Limitaciones}
\begin{itemize}
\item Se usó una única partición de prueba (503 y 98 casos). Con 15 diabéticos en
prueba, un solo paciente cambia el recall en 6.7 puntos. No se calcularon intervalos de
confianza ni pruebas de significancia.
\item El catálogo tiene tres variantes por familia y cada una cambia un solo
hiperparámetro. Es deliberado, para poder atribuir los efectos, pero no explora
combinaciones (por ejemplo, más árboles y hojas mínimas a la vez).
\item Todos los modelos usan el umbral de probabilidad 0.5. Ajustarlo con validación
cruzada podría aumentar el recall de la clase minoritaria.
\item La clase de Diabetes se construyó con el criterio histórico $\text{glyhb}>7$, y la
versión de Potabilidad usada son los casos completos del archivo público. Los resultados
son educativos: no constituyen un diagnóstico médico ni una certificación de potabilidad.
\item La importancia por permutación mide de qué depende el modelo, no relaciones
causales.
\end{itemize}
